\documentclass[10pt,a4paper]{amsart}
\usepackage[margin=19mm]{geometry}
\usepackage{amsmath,amssymb,mathtools}
\usepackage[scr=rsfs,cal=euler]{mathalfa}
\usepackage{xfrac}
\usepackage[unicode,hidelinks,bookmarks=false]{hyperref}
\newcommand{\ie}{i.e.\ }
\newcommand{\cf}{cf.\ }
\newcommand{\resp}{respectively\ }
\newcommand{\Subsection}{\subsection}
\providecommand{\nobreakdash}{\nobreak\hbox{-}}
\setlength{\emergencystretch}{2em}
\multlinegap0pt
\usepackage{amssymb}
\usepackage{mathtools}
\usepackage{multirow}
\usepackage{nicefrac}
\usepackage{enumerate}
\newtheorem{lem}{Lemma}
\title{Ciliberto--Di Gennaro conjecture for sextic hypersurfaces}
\author{Ksenia Kvitko}
\date{Source: arXiv:2505.06742v2 (CC BY 4.0)}
\begin{document}
\maketitle

\noindent\emph{Source and licence.} Ciliberto--Di Gennaro conjecture for sextic hypersurfaces. Version arXiv:2505.06742v2 (CC BY 4.0), distributed by arXiv under the Creative Commons Attribution 4.0 International licence, http://creativecommons.org/licenses/by/4.0/, as recorded in the arXiv licence metadata for this exact version and shown on the abstract page https://arxiv.org/abs/2505.06742v2. Full licence notice: \texttt{licenses/arxiv-2505.06742-CC-BY-4.0.txt}.\par\smallskip

\noindent\textit{Editorial scope.} Counted unit: the article's lem auxiliary\_cii (source0.tex lines 412--428). The article's complete original proof of it is delivered with it (source0.tex lines 430--463, one \texttt{proof} environment of 34 lines). No mathematics is written, repaired or re-derived: the statement and the proof are the article's own lines, in the article's own notation, and no retained line is substituted or re-flowed.  No further source-local context is needed: every cross-reference of every retained line resolves inside the two delivered spans.  Nothing was omitted from any retained span, so the package carries no omission record.  No retained line contains a citation, so the wrapper carries no bibliography and no bibliography entry.  No retained line contains a figure environment, an image-inclusion command or drawing code, so no image is delivered and the package declares no asset dependency.  Editorial formatting only, in addition: an AMS \texttt{amsart} preamble that restates the article's own declarations and macros used by the retained lines, namely the environments \texttt{lem} on separate counters, together with the standard packages those lines need.  The retained environments print in this document as Lemma 1 in order of appearance, whereas the article prints the same items with the numbering of its own full document.  The complete native source of the article as distributed in the arXiv e-print for this exact version is bundled unchanged as \texttt{sources/177903/source0.tex}.
\begin{lem}\label{auxiliary_cii}
		Given an artinian Gorenstein ring $S/I$ of socle degree $2d-4$, $d\geqslant 6$, assume the following conditions are satisfied:
		\begin{enumerate}
		\item[\emph{(1)}] $h_I(1)\geqslant 3$;
		\item[\emph{(2)}] $h_I(k)=h_I(2d-4-k)\geqslant 2k+2$ for $2\leqslant k \leqslant d-4$;
		\item[\emph{(3)}]  $h_I(k)\geqslant 2d-6$ for $d-3\leqslant k \leqslant d-1$.
		\end{enumerate}
		Then the inequality holds
		\begin{equation}\label{inequality_for_singularities}\tag{4.2}
			 \sum\limits_{k=0}^{2d-4} h_I(k)>2(d-2)(d-1)
		\end{equation}
		except when $d= 7$ and $I$ has $h$-vector
		$$(1,3,6,8,8,8,8,8,6,3,1)$$
		or when $d=6$ and $h$-vector of $I$ is
		$$(1,3,6,6,6,6,6,3,1),\hspace{10pt} (1,3,6,6,7,6,6,3,1), \hspace{10pt}(1,3,6,7,6,7,6,3,1), \hspace{10pt} (1,3,7,6,6,6,7,3,1), $$
		 $$(1,4,6,6,6,6,6,4,1), \hspace{10pt} (1,3,6,6,8,6,6,3,1).$$
\end{lem}

\begin{proof}
	In order to prove inequality \eqref{inequality_for_singularities}, consider an auxiliary number sequence $h'(k)$, $k=0,\dots 2d-4$. Let $h'(0)=1$, $h'(1)=3$, $h'(k)=2k+1$ for $2\leqslant k \leqslant d-3$, and $h'(d-2)=2(d-2)$.
	

	Using conditions (1)-(3), we compare the values $h_I(k)$ and $h'(k)$ in Table \ref{table}.
	
	\begin{table}[h!]
	\centering
		\begin{tabular}{| c | c | c | c | c | c | c | c | c | c | c |}\hline 
		$k$ & 0 & 1 & 2  & \dots & $m$ & \dots & $d-4$ & $d-3$ & $d-2$ & \dots \\ \hline
		$h_I(k)$ & 1 & $\geqslant 3$ & $\geqslant 6$ & \dots & $\geqslant 2m+2$ & \dots & $\geqslant 2d-6$ &  $\geqslant 2d-6$ & $\geqslant 2d-6$ & \dots   \\ \hline
		$h'(k)$ & 1 & 3 & 5 & \dots & $2m+1$ & \dots & $2d-7$ & $2d-5$ & $2d-4$ & \dots \\ \hline
	\end{tabular}
	\caption{Values of Hilbert function of ideal $I$.}
	\label{table}
	\end{table}
	One can see that for fixed $d$ the sum of $h_I(k)$, $0\leqslant k\leqslant 2d-4$, is at least $2(d^2-2d-5)$ while the~sum of $h'(k)$ equals exactly $2(d-2)(d-1)$. So, one has 
		\begin{multline*}
		\sum\limits_{k=0}^{2d-4}\bigg(h_I(k)-h'(k)\bigg)\geqslant  \\
		\geqslant 2\sum\limits_{k=2}^{d-4}\bigg((2k+2)-(2k-1)\bigg)+ 2 \bigg( (2d-6) -(2d-5)\bigg) +\bigg((2d-6)-2(d-2)\bigg)=\\
		=2d-14,
	\end{multline*}
	which proves the lemma for $d>7$. 
	
	Now, suppose $d=7$. Then $2(d^2-2d-5)=2(d-2)(d-1)=60$ and therefore
	$$\sum\limits_{k=0}^{2d-4}h_I(k)\geqslant 1+3+6+8+\dots+8+6+3+1=60= 2(d-2)(d-1).$$ Here, the equality is attained if $I$ has $h$-vector $(1,3,6,8,8,8,8,8,6,3,1)$. 
	
	Let $d=6$. In this case, we have 
	$$\sum\limits_{k=0}^{2d-4}h_I(k)\geqslant 1+3+6+6+6+6+6+3+1= 38$$
	and $2(d-2)(d-1)=40$. Denote the sum by $S$. If $S=38$, then $h$-vector is $(1,3,6,6,6,6,6,3,1)$. \mbox{If $S=39$,} then due to the Gorenstein duality the middle value $h_I(d-2)$ is greater than the lower bound by 1. So, the $h$-vector is $(1,3,6,6,7,6,6,3,1)$. If $S=40$, then we have four possibilities to distribute additional 2 among values $h_I(k)=h_I(2d-4-k)$ and therefore $h$-vector is one of the following
	$$(1,4,6,6,6,6,6,4,1), \hspace{10pt}(1,3,7,6,6,6,7,3,1), \hspace{10pt} (1,3,6,7,6,7,6,3,1), \hspace{10pt} (1,3,6,6,8,6,6,3,1).$$
	For $S>40$ the inequality \eqref{inequality_for_singularities} holds.

\end{proof}

\end{document}
